\bookchapter{Anhang F: Faltbildschirm-Kauf und -Nutzung}

\chapterintro{Wichtige Fortschritte 2024–2025: Scharnierlebensdauer meist mit 300.000 Klappzyklen spezifiziert (bei 100-mal täglichem Aufklappen ca. 8 Jahre), UTG-Dünnglas ersetzt die frühere Kunststofffolie, die Knickfalte ist sichtbar, aber das Gefühl beim Anfassen bereits sehr dezent, Preise von der 10.000-Yuan-Klasse auf 5.000–6.000 Yuan gesunken.}

\booksection{Erstens. Wie weit sind Klappbildschirme heute entwickelt?}

Wenn dir ein starres Smartphone völlig ausreicht und Sie keinen Schmerzpunkt “Bildschirm zu klein” haben, kaufen Sie keinen Klappbildschirm. Er ist ein “i-Tüpfelchen”-Gerät, kein “Retter in der Not”.\textbf{Zweitens. Für wen eignet sich ein Klappbildschirm?}Laut echtem Nutzerfeedback haben folgende Personengruppen eine hohe Wahrscheinlichkeit, den Klappbildschirm “toll” zu finden:

Intensive Leser/Dokumentenverarbeiter

\begin{itemize}
  \item \textbf{: Wer oft PDFs am Smartphone liest, Dokumente kommentiert oder lange Texte liest. Aufgeklappt bietet der Innenbildschirm rund 7,5 Zoll, A4-PDFs lassen sich ohne Zoom lesen – ein Qualitätssprung.}Geschäftsreisende/Besprechungs-Profis
  \item \textbf{: Wer kein Tablet mitnehmen will, aber das Smartphone als zu klein empfindet. Ein Horizontal-Klappbildschirm zeigt PPTs, Verträge und E-Mails im Zweifenster-Vergleich deutlich effizienter als ein starres Gerät.}"Aufsteller-Kreative"
\end{itemize}

Wichtige Fortschritte 2024–2026: Scharnierlebensdauer allgemein mit 500.000 Faltzyklen angegeben (bei 100-mal täglich 13 Jahre), UTG-Ultradünglas ersetzt frühe Kunststofffolien, Falz mit bloßem Auge sichtbar aber kaum fühlbar, Preise sinken von über 10.000 auf 4.000–5.000 Yuan.

\begin{tcbinfo}{Tragbarkeits- und Großbildschirm-Fans}
: Vertikal-Klappbildschirme passen zugeklappt in kleine Taschen, sind aufgeklappt normale Smartphones und können den Außenbildschirm als Spiegel oder für Benachrichtigungen nutzen.

\end{tcbinfo}

\booksection{Zweitens. Für wen eignet sich ein Klappbildschirm?}

Wer das Smartphone nur zum Kurzvideo-Scrollen, WeChat-Chatten und Telefonieren nutzt – der Klappbildschirm bringt keinen Vorteil, ist dafür schwerer und dicker.

\begin{enumerate}
  \item \textbf{Personen mit kleinen Händen oder Sehnenscheidenentzündung – Horizontal-Klappbildschirme wiegen aufgeklappt meist 230 g+, einhändiges Halten ist anstrengend.}Bei knappem Budget oder wenn ein Gerät länger als drei Jahre halten soll – Klappbildschirme haben einen niedrigeren Gebrauchtwert als starre Geräte und höhere Reparaturkosten.
  \item \textbf{Drittens. Kaufberatung: Achte auf diese Kennzahlen}Lass dich nicht vom Wort "Klapp" blenden. Achte beim Kauf auf:
  \item \textbf{Scharnier-Art}: Wassertropfen-Scharnier (flache Knickfalte, lange Lebensdauer, heute Standard) ist besser als U-Scharnier (tiefe Knickfalte, frühe Lösung). Beim Verkäufer nachfragen oder Testberichte prüfen.
  \item \textbf{Innenbildschirm-Material}: Muss UTG-Dünnglas sein, nicht die frühere Kunststofffolie. UTG fühlt sich fast wie Glas an; bei Kunststofffolie hinterlässt der Fingernagel Eindrücke und Kratzer.
\end{enumerate}

\begin{tcbwarn}{Knickfalten-Erscheinung}
\begin{itemize}
  \item : Im Geschäft selbst aufklappen und ansehen. Akzeptables Niveau 2025: Bei eingeschaltetem Bildschirm frontal kaum sichtbar, im schrägen Licht eine flache Furche erkennbar, beim Darüberstreichen leicht spürbar, aber nicht scharfkantig.
  \item Außenbildschirm-Größe
  \item : Bei Horizontal-Klappbildschirmen gilt "je größer, desto besser" (mindestens 6,4 Zoll), damit er zugeklappt wie ein normales Smartphone nutzbar ist. Bei Vertikal-Klappbildschirmen ist ein größerer Außenbildschirm praktischer (ab 4 Zoll für Nachrichten und Navigation).
\end{itemize}

\end{tcbwarn}

\booksection{Drittens. Kaufberatung: Achte auf diese Kennzahlen}

: Horizontal-Klappbildschirme wiegen aufgeklappt meist 230–260 g; unter 230 g gilt als gut. Zugeklappt 12–14 mm; unter 12 mm gilt als dünn. Vertikal-Klappbildschirme sollten zugeklappt bequem in die Hemdtasche passen.

\begin{itemize}
  \item \textbf{Service und Garantie}: Vor allem die Reparaturkosten für einen gebrochenen Innenbildschirm prüfen. Manche Marken bieten im ersten Jahr einen kostenlosen Austausch des Innenbildschirms an, bei anderen kostet er 2.000–4.000 Yuan aus eigener Tasche – das ist die größte versteckte Kostenstelle des Klappbildschirms.
  \item \textbf{Checkliste für den Store-Test}10-mal selbst auf- und zuklappen, prüfen, ob das Scharnier sanft geht und ob es Geräusche gibt.
  \item \textbf{Aufgeklappt schräg ins Licht halten, um die Knickfalte zu sehen; dann frontal bei eingeschaltetem Bildschirm prüfen.}: Gehen Sie in ein Geschäft und klappen Sie es selbst auf. Der 2026er Standard: Bei eingeschaltetem Bildschirm vorne kaum sichtbar, bei schrägem Licht eine leichte Rinne erkennbar, beim Darüberstreichen eine leichte Vertiefung, aber nicht scharfkantig.
  \item \textbf{Die Aufsteller-Position (90 Grad) ausprobieren: Fotos und Videos testen.}Repaurpreis für den Innenbildschirm und Garantiebestimmungen erfragen.
  \item \textbf{Viertens. Spezielle Nutzungstechniken des Klappbildschirms}Nach dem Kauf bieten diese Anwendungen ein Erlebnis, das starre Geräte nicht bieten.
  \item \textbf{1. Split-Screen und Multitasking}Bei Horizontal-Klappbildschirmen unterstützen die meisten Systeme nach dem Aufklappen einen oberen/unteren oder linken/rechten Split-Screen. Typische Szenarien: oben Online-Kurs-Video, unten Notizen; links WeChat, rechts Dokument; oben Karten-Navigation, unten Musik. Praktischer als bei starren Geräten, weil der Bildschirm groß genug ist und jede Hälfte nicht beengt wirkt.
\end{itemize}

\begin{tcbpractice}{2. Aufsteller-Modus (90 Grad gefaltet)}
\begin{enumerate}
  \item Das ist das "Ass" des Klappbildschirms.
  \item Selfie/Gruppenfoto
  \item : Auf den Tisch gestellt, 90 Grad, oberer Bildschirm als Sucher, unterer als Auslöser – niemand muss um Hilfe bitten.
  \item Zeitraffer/Nachtaufnahme
  \item : Ebenfalls stabil auf dem Tisch, ein Stativ entfällt.
\end{enumerate}

\end{tcbpractice}

\booksection{Viertens. Spezielle Nutzungstechniken des Klappbildschirms}

: Hände frei, der Gegenüber sieht dich in einem natürlicheren Winkel.

\booksubsection{Video schauen}

: Als kleiner Fernseher auf den Tisch gestellt – besonders praktisch auf der Klapptischchen in Flugzeug oder Zug.

\booksubsection{Außenbildschirm-Vorschau + Innenbildschirm-Aufnahme}

: Wenn du mit der Rückkamera dich selbst filmst, zeigt der Außenbildschirm das Bild live – bessere Bildkomposition.

\begin{itemize}
  \item \textbf{3. Außenbildschirm-Schnellblick}Bei zugeklappten Vertikal-Klappbildschirmen kann der Außenbildschirm: Uhrzeit, Benachrichtigungen, Wetter anzeigen, WeChat-Nachrichten beantworten (Sprache-zu-Text oder Schnellantwort), Transitkarte nutzen, Navigation anzeigen, als Spiegel dienen (über die Frontkamera). Nicht jedes Mal aufklappen – stromsparend und bequem.
  \item \textbf{4. Computer-Modus}Einige Horizontal-Klappbildschirme (Huawei, Samsung) können nach dem Aufklappen in einen "Computer-Modus" wechseln; mit externem Display, Tastatur und Maus ähnelt er einem einfachen PC und reicht für temporäre Dokumentenbearbeitung oder E-Mail-Beantwortung. Kein primäres Produktivitätstool, aber in Notsituationen sehr willkommen.
  \item \textbf{Fünftens. Pflegehinweise: Den Klappbildschirm länger am Leben halten}Am empfindlichsten sind der Innenbildschirm und das Scharnier. Diese Gewohnheiten verlängern die Lebensdauer deutlich:
  \item \textbf{Bei Kälte nicht gewaltsam aufklappen}: Unter 5 °C wird UTG-Glas spröde; vor dem Aufklappen das Gerät einige Minuten in der Tasche wärmen. Das steht in jedem offiziellen Handbuch, wird aber ignoriert.
  \item \textbf{Fingernägel/Hartgegenstände vom Innenbildschirm fernhalten}: Der Innenbildschirm hat zwar eine Schutzschicht, ist aber weicher als der Außenbildschirm. Mit dem Fingerballen bedienen, nicht mit dem Fingernagel tippen.
\end{itemize}

\booksubsection{Vor dem Zuklappen sicherstellen, dass sich nichts auf dem Bildschirm befindet}

: Ein Sandkorn, das beim Zuklappen dazwischen gerät, kann den Innenbildschirm direkt beschädigen. Vor dem Zuklappen kurz hinschauen.

\booksubsection{Keine normalen Panzerglas-Folgen kleben}

: Der Innenbildschirm darf nur mit der offiziellen Spezialfolie oder einer UTG-Schutzfolie versehen werden; normale Panzerglas-Folgen beeinträchtigen Klappfunktion und Touch. Folie beim offiziellen Service anbringen lassen.

\booksection{Fünftens. Pflegehinweise: Den Klappbildschirm länger am Leben halten}

: Zwar ist die Mehrzahl mit IPX8 als wasserdicht ausgewiesen, aber der Staubschutz ist meist gering. Am Strand, auf Baustellen und in ähnlichen Umgebungen vorsichtig sein.

\begin{itemize}
  \item \textbf{Auf- und Zuklappen mit gleichmäßiger Kraft}: Nicht einhändig aufreißen; mit beiden Händen von der Mitte aus zügig und gleichmäßig aufklappen ist am stabilsten.
  \item \textbf{Häufige Innenbildschirm-Schadensursachen}Beim Zuklappen werden Smartphone-Halterung oder Ladekabel eingeklemmt und der Innenbildschirm reißt direkt.
  \item \textbf{In der Tasche zusammen mit Schlüsseln aufbewahrt – der Außenbildschirm überlebt, aber der Innenbildschirm bekommt Risse durch die Schlüssel.}Im Winter aus der Kälte in einen warmen Raum bringen und sofort auf- bzw. zuklappen – Temperaturschock verändert die Glasspannung.
  \item \textbf{Nach einem Jahr Nutzung den Eindruck haben, die Knickfalte sei tiefer geworden, und die Werksfolie abziehen – das Abziehen kann den Innenbildschirm mit beschädigen.}Sechstens. Preis- und psychologische Vorbereitung des Klappbildschirms
  \item \textbf{Vor dem Kauf solltest du folgende Rechnungen im Kopf haben:}Die Knickfalte bleibt
  \item \textbf{: 2025 kann die Technik "vollständig faltenfrei" nicht leisten, sondern nur "bei eingeschaltetem Bildschirm frontal kaum sichtbar". Wer Perfektionist ist und täglich die Knickfalte stört, sollte nicht kaufen.}Schwerer und dicker
\end{itemize}

\textbf{: Horizontal-Klappbildschirme wiegen aufgeklappt 230 g+, einhändiges Halten ist anstrengender als beim starren Gerät; zugeklappt 12–14 mm dick – die Hosentasche wird dicker.}\par

\begin{itemize}
  \item Akkulaufzeit meist kürzer
  \item : Wegen des Scharnier-Mechanismus wird der Akkuraum komprimiert; in derselben Preisklasse ist die Akkulaufzeit meist schlechter als bei starren Geräten.
  \item Wertverlust
  \item : Nach einem Jahr liegt der Gebrauchtpreis meist bei 50–60 \% des Neupreises – schneller als bei starren Geräten.
\end{itemize}

\booksection{Sechstens. Preis- und psychologische Vorbereitung des Klappbildschirms}

: Innenbildschirm-Reparatur 2.000–4.000 Yuan, Scharnier-Reparatur 500–1.500 Yuan. Vor dem Kauf Garantiebestimmungen prüfen.

\begin{itemize}
  \item \textbf{App-Anpassung uneinheitlich}: Die 2026er Technik kann nicht "komplett spurlos", nur "bei eingeschaltetem Bildschirm vorne kaum sichtbar". Wenn Sie Perfektionist sind und die Falz ständig stört – nicht kaufen.
  \item \textbf{Der Klappbildschirm ist kein "besseres Smartphone", sondern "ein Smartphone in anderer Form". Kauf ihn, weil du seine einzigartige Form brauchst, nicht weil er fortschrittlicher ist.}– Zur Kaufphilosophie des Klappbildschirms
  \item \textbf{Siebtens. Auswahl der Marken (Stand 2025)}Keine konkreten Modelle (die veralten), nur Marken-Charakteristika, damit du dich schnell orientieren kannst:
  \item \textbf{Huawei Mate X-Serie}: Branchenbenchmark für Horizontal-Klappbildschirme, starke HarmonyOS-Ökosystem-Synergie, Außenbildschirm-Proportionen nahe einem normalen Smartphone, zugeklappt als Hauptgerät nutzbar. Hoher Preis, aber relativ bester Wertverlust.
  \item \textbf{Samsung Galaxy Z Fold/Flip}: Klappbildschirm-Pionier, Scharnier ausgereift, Z Flip ist der Vertikal-Klapp-Vertreter. Im Inland ist das After-Sales-Netz kleiner, Ersatzteile sind teuer.
  \item \textbf{Xiaomi MIX Fold/Flip}: Preis-Leistungs-Route, großzügige Ausstattung, relativ günstiger Preis. Geeignet für Nutzer, die Neues ausprobieren möchten, ohne zu viel auszugeben.
\end{itemize}

vivo X Fold Serie

: Unter den horizontal faltbaren Geräten hat das Verhältnis von Außen- zu Innenbildschirm die größte Ähnlichkeit mit einem normalen Smartphone; zugeklappt als Barren-Handy bietet es ein gutes Nutzungserlebnis.

\booksection{VII. Welche Marke wählen (2026)}

: Unter den horizontal faltbaren Geräten eher kompakt, einhändig bedienbar, geeignet für Nutzer, denen große Faltbildschirme zu schwer sind.

\begin{itemize}
  \item \textbf{Huawei Pocket Serie}: Unter den vertikal faltbaren Geräten der größte und praktischste Außenbildschirm; Kosmetikspiegel- und Sonnenschutz-Erkennungsfunktionen für weibliche Nutzer sind detailliert umgesetzt.
  \item \textbf{Kaufempfehlungen: Bei ausreichendem Budget und Fokus auf das Ökosystem wähle Huawei Mate X; für den ausgereiftesten Scharniermechanismus wähle Samsung Z Fold; für das beste Preis-Leistungs-Verhältnis wähle Xiaomi/vivo; bei vertikalen Faltbildschirmen vergleiche Huawei Pocket und Samsung Z Flip.}Letzter Hinweis
  \item \textbf{Die Faltbildschirm-Technologie entwickelt sich jedes Jahr weiter; dieser Anhang wurde 2025 verfasst. Vor dem Kauf unbedingt die neuesten Tests und Nutzerfeedbacks prüfen, mit Fokus auf: gemessene Scharnierlebensdauer, Falzverhalten und Reparaturrichtlinien für den Innenbildschirm. Ein persönlicher Test im Geschäft ist immer wichtiger als Online-Bewertungen.}Zuletzt gelesen bei: ...
  \item \textbf{Weiterlesen}Ignorieren
  \item \textbf{OPPO Find N-Serie}:Kompaktere Modelle bei horizontalen Klapphandys, Einhaendig bedienbar, Fuer Nutzer, denen grosse Klapphandys zu schwer sind.
  \item \textbf{HuaweiPocket-Serie}:Bei vertikalen Klapphandys die groesste und praktischste Aussenanzeige, Kosmetikspiegel, Sonnenschutz-Erkennung und weitere frauenorientierte Funktionen sind detailliert umgesetzt.
\end{itemize}

Kaufempfehlungen: Bei ausreichendem Budget und Fokus auf Oekosystem: HuaweiMate X; Bei Wunsch nach ausgereiftem Scharnier: SamsungZ Fold; Bestes Preis-Leistungs-Verhaeltnis: Xiaomi/vivo; Bei vertikalem Klapphandy: HuaweiPocketund SamsungZ Flip.

\begin{tcbinfo}{Letzte Erinnerung}
Faltbare Displays entwickeln sich jedes Jahr weiter. Dieser Anhang wurde 2026 verfasst. Vor dem Kauf unbedingt aktuelle Testberichte und Nutzer-Feedback prüfen – Schwerpunkte: Scharnierlebensdauer im Test, Falz-Darstellung, Reparaturrichtlinie für Innenbildschirm. Selber Anfassen im Laden ist immer wichtiger als Testberichte.

\end{tcbinfo}
